\subsection{From a black-box judgment to a gradient}

An RL rollout produces a response and one scalar judgment. The checker supplies no
corrected tokens and no derivative, so the model can only reweight its own sampled
behavior: make successful decisions more probable and unsuccessful decisions less
probable. Differentiating expected reward gives the policy-gradient identity
\begin{equation}
  \nabla_\theta J
  =\mathbb{E}_{o\sim\pi_\theta}
    \left[A(o)\nabla_\theta\log\pi_\theta(o\mid q)\right].
  \label{eq:compact-policy-gradient}
\end{equation}
The log-probability is not an arbitrary loss. It is the sample-based derivative of ``put
more probability on behavior with positive advantage.'' For an autoregressive response,
$\log\pi_\theta(o\mid q)$ is the sum of its conditional token log-probabilities, so one
outcome advantage is broadcast across the sampled token decisions.

Raw reward is a noisy direction because its meaning depends on the prompt. Advantage
subtracts the reward expected in the same situation: an $8$ is good when the policy
usually scores $2$ and bad when it usually scores $9$. PPO commonly estimates this
baseline with a learned value function. GRPO instead samples $G$ sibling responses to one
prompt and uses their reward statistics:
\begin{equation}
  A_i=\frac{r_i-\mu_q}{\sigma_q+\epsilon_A},
  \qquad
  \mu_q=\frac{1}{G}\sum_j r_j,
  \qquad
  \sigma_q^2=\frac{1}{G}\sum_j(r_j-\mu_q)^2.
  \label{eq:compact-grpo-advantage}
\end{equation}
The group forms a small tournament: above-average attempts are reinforced and
below-average attempts are suppressed. This avoids training a critic but pays for several
rollouts per prompt and retains Monte Carlo noise. Division by the group standard deviation
equalizes reward scale across prompts; it is a weighting heuristic rather than a necessary
property of baseline subtraction.
